\documentclass[11pt,a4paper]{article}

% ------------------------------------------------------------
% Packages
% ------------------------------------------------------------
\usepackage[utf8]{inputenc}
\usepackage[T1]{fontenc}
\usepackage[brazilian]{babel}
\usepackage{lmodern}
\usepackage{amsmath,amssymb}
\usepackage{geometry}
\usepackage{setspace}
\usepackage{hyperref}

% ------------------------------------------------------------
% Page setup
% ------------------------------------------------------------
\geometry{
    left=2.5cm,
    right=2.5cm,
    top=2.5cm,
    bottom=2.5cm
}

\onehalfspacing

% ------------------------------------------------------------
% Title
% ------------------------------------------------------------
\title{Nota sobre Emiss\~oes Brutas e L\'iquidas na Constru\c{c}\~ao dos Choques de D\'ivida}
\author{}
\date{\today}

\begin{document}

\maketitle

\section*{Nota sobre emiss\~oes brutas e l\'iquidas}

A emiss\~ao bruta de mercado por classe de d\'ivida \(j\in\{N,F,R\}\) pode ser definida diretamente a partir dos leil\~oes prim\'arios como
\[
G_t^j
=
Venda_t^j
+
ExtraVenda_t^j,
\]
onde \(N\) corresponde a LTN e NTN-F, \(F\) a LFT, e \(R\) a NTN-B. Esta medida captura o valor financeiro colocado pelo Tesouro no mercado em cada classe de t\'itulo.

No arquivo de leil\~oes, \(Venda_t^j\) representa as vendas ordin\'arias de t\'itulos da classe \(j\), enquanto \(ExtraVenda_t^j\) representa vendas extraordin\'arias da mesma classe. Assim, \(ExtraVenda_t^j\) n\~ao \'e recompra, resgate ou troca; trata-se de uma coloca\c{c}\~ao adicional de t\'itulos no mercado, possivelmente realizada fora do leil\~ao ordin\'ario regular ou em resposta a condi\c{c}\~oes espec\'ificas de caixa, demanda, liquidez ou gest\~ao da d\'ivida. Por ter o mesmo sinal econ\^omico de uma venda ordin\'aria, ela deve ser somada \`a emiss\~ao bruta.

A emiss\~ao l\'iquida, por sua vez, n\~ao \'e uma medida única. Uma defini\c{c}\~ao operacional baseada em leil\~oes seria
\[
I_t^{j,\mathrm{liq,mercado}}
=
Venda_t^j
+
ExtraVenda_t^j
-
Compra_t^j
-
ExtraCompra_t^j
+
TrocaIn_t^j
-
TrocaOut_t^j,
\]
caso seja poss\'ivel identificar separadamente os dois lados das opera\c{c}\~oes de troca. Nesta express\~ao, \(Compra_t^j\) representa compras ou recompras ordin\'arias de t\'itulos da classe \(j\) pelo Tesouro, enquanto \(ExtraCompra_t^j\) representa compras extraordin\'arias. Ambas t\^em sinal negativo porque reduzem o volume l\'iquido de t\'itulos da classe \(j\) em poder do mercado. Em outras palavras, \(ExtraCompra_t^j\) \'e uma recompra extraordin\'aria ou opera\c{c}\~ao extraordin\'aria de retirada de t\'itulos do mercado, e n\~ao uma emiss\~ao.

As vari\'aveis \(TrocaIn_t^j\) e \(TrocaOut_t^j\) capturam os dois lados de opera\c{c}\~oes de troca. Em uma troca, o Tesouro entrega ou emite um t\'itulo novo e recebe de volta um t\'itulo antigo. Assim, \(TrocaIn_t^j\) denota o valor financeiro dos t\'itulos da classe \(j\) entregues pelo Tesouro ao mercado na troca, isto \'e, o lado de emiss\~ao da opera\c{c}\~ao. J\'a \(TrocaOut_t^j\) denota o valor financeiro dos t\'itulos da classe \(j\) recebidos pelo Tesouro e retirados do mercado na troca, isto \'e, o lado de resgate da opera\c{c}\~ao. Uma mesma opera\c{c}\~ao de troca pode, portanto, aumentar a emiss\~ao l\'iquida de uma classe e reduzir a de outra. Sem a separa\c{c}\~ao entre \(TrocaIn_t^j\) e \(TrocaOut_t^j\), as trocas podem misturar emiss\~ao de um t\'itulo novo com resgate de outro t\'itulo, possivelmente de classe diferente.

J\'a a emiss\~ao l\'iquida oficial agregada da DPF ou da DPMFi \'e usualmente definida como
\[
I_t^{\mathrm{liq,oficial}}
=
Emissoes_t
-
Resgates_t,
\]
em que os resgates podem incluir vencimentos de principal, pagamentos de cupons, recompras, compras antecipadas, opera\c{c}\~oes de troca, cancelamentos e outras opera\c{c}\~oes diretas, dependendo da cobertura da estat\'istica utilizada. Finalmente, uma medida inferida a partir da varia\c{c}\~ao do estoque satisfaz uma identidade cont\'abil do tipo
\[
\Delta B_t^j
=
I_t^{j,\mathrm{liq}}
+
Juros_t^j
+
Indexacao_t^j
+
Ajustes_t^j,
\]
ou, de forma equivalente,
\[
I_t^{j,\mathrm{liq}}
=
\Delta B_t^j
-
Juros_t^j
-
Indexacao_t^j
-
Ajustes_t^j.
\]
Esta \'ultima constru\c{c}\~ao \'e particularmente delicada para LFT e NTN-B, pois a varia\c{c}\~ao do estoque incorpora mecanicamente a apropria\c{c}\~ao da Selic e a corre\c{c}\~ao pelo IPCA, respectivamente.

Por esses motivos, a emiss\~ao l\'iquida tende a misturar v\'arias margens de decis\~ao e de ajuste: a decis\~ao de coloca\c{c}\~ao prim\'aria, o calend\'ario de vencimentos, recompras, opera\c{c}\~oes de troca, pagamentos de cupons, indexa\c{c}\~ao, juros apropriados e ajustes cont\'abeis. Al\'em disso, algumas dessas margens s\~ao respostas end\'ogenas ao ambiente macroecon\^omico e financeiro. Em per\'iodos de estresse, por exemplo, o Tesouro pode alterar recompras, trocas e composi\c{c}\~ao das coloca\c{c}\~oes em resposta a condi\c{c}\~oes de liquidez, risco soberano, inclina\c{c}\~ao da curva de juros ou demanda dos investidores.

Assim, para a especifica\c{c}\~ao principal, parece prefer\'ivel construir os choques a partir das emiss\~oes brutas por instrumento. A emiss\~ao bruta \'e mais pr\'oxima da decis\~ao marginal de coloca\c{c}\~ao do Tesouro e permite separar, na primeira etapa, a parte previs\'ivel da emiss\~ao associada a vencimentos programados, composi\c{c}\~ao defasada da d\'ivida, PAF, reserva de liquidez e condi\c{c}\~oes macrofinanceiras:
\[
g_t^j
=
\alpha^j
+
\delta_0^j m_{t\mid t-1}^j
+
\delta_1^j m_{t+1\mid t-1}^j
+
\Lambda^{j\prime}Z_{t-1}
+
\mu_{m(t)}^j
+
\nu_t^j.
\]
O res\'iduo \(\nu_t^j\) pode ent\~ao ser interpretado como uma inova\c{c}\~ao na coloca\c{c}\~ao bruta da classe \(j\), condicional ao conjunto de informa\c{c}\~ao observado. Essa abordagem facilita a distin\c{c}\~ao entre a decis\~ao de emitir e outras decis\~oes de gest\~ao da d\'ivida que o Tesouro pode tomar simultaneamente em resposta ao ambiente macrofinanceiro.

As emiss\~oes l\'iquidas continuam sendo úteis como exerc\'icio de robustez, desde que a defini\c{c}\~ao utilizada seja explicitada com precis\~ao. Em particular, deve-se distinguir entre emiss\~ao l\'iquida de mercado baseada em leil\~oes, emiss\~ao l\'iquida oficial agregada e emiss\~ao l\'iquida inferida por varia\c{c}\~oes de estoque. No entanto, por transpar\^encia e menor risco de contaminar os choques com juros, indexa\c{c}\~ao e ajustes end\'ogenos, a identifica\c{c}\~ao principal deve se basear nas emiss\~oes brutas por instrumento.

\end{document}
